\par
\SourceLeaf{368}%
\pdfbookmark[0]{Anzeigen. Kries: Mathematische Geographie. 1814.}{anzeigen-kries-1814}
\begin{center}
\rule{.85\textwidth}{.3pt}\par
Göttingische gelehrte Anzeigen.\quad 1814 Juni 13.\par
\rule{.85\textwidth}{.3pt}
\end{center}
\hspace*{\parindent}\textit{Lehrbuch der mathematischen Geographie von} \textsc{Friedrich Kries}, Professor am Gymnasium zu Gotha. Mit sieben Kupfertafeln. 236 Seiten in Octav. Leipzig, bei \textsc{G. J. Göschen}.

Der Plan des Verfassers bei Abfassung dieses Lehrbuchs für eine Wissenschaft, welche für jeden Gebildeten ein so vielseitiges Interesse hat, ging dahin, zwischen den dürftigen und oberflächlichen Abrissen derselben, die den Lehrbüchern der politischen Erdbeschreibung vorangeschickt zu werden pflegen, und sich nur auf die Aufzählung von Hauptresultaten beschränken, ohne sie durch mathematische Behandlung zu begründen oder zu erläutern, --- und den grössern Werken, welche feinere, weniger allgemein verbreitete Kenntnisse der höhern Mathematik voraussetzen, eine schickliche Mittelstrasse zu treffen. In einem solchen Werke erwartet man nicht neue Aufklärungen, die die Wissenschaft selbst weiter bringen, sondern nur, dass eine zweckmässige Auswahl aus dem Bekannten mit Ordnung, Gründlichkeit und Klarheit dargestellt werde, und dieses Ziel hat der Verf. in der That erreicht. Er handelt in zehn Abschnitten von der Gestalt des Erdkörpers im Allgemeinen; von der mathematischen Eintheilung der Erdkugel und ihrer Grösse; von der Umdrehung derselben um ihre Axe und den damit zusammenhängenden Erscheinungen; von den Mitteln, die geographische Breite eines Orts zu bestimmen, und eine Mittagslinie zu ziehen; von der Bewegung der Erde um die Sonne; von der Eintheilung der Himmels- und der Erdkugel in Beziehung auf die Bewegung der Erde um die Sonne, und den Erscheinungen, die auf der Erde aus dieser Bewegung entstehen; von der Zeitbestimmung und den Mitteln zur Bestimmung der geographischen Länge; von der sphäroidischen Gestalt der Erde; von der Verfertigung künstlicher Erdkugeln und der Landkarten; vom Gebrauch der künstlichen Erdkugel zur Auflösung mathematisch geographischer Aufgaben. Wir können nicht anders, als dieser Anordnung und Auswahl im Allgemeinen unsern Beifall geben, wenn gleich unsrer Ansicht nach hie und da noch einige Gegenstände, die nicht berührt sind, hätten aufgenommen, und dagegen andere z.\,B. die verschiedenen Projectionsarten der Karten allenfalls etwas kürzer hätten abgehandelt werden können. So hätten wir unter andern einige Anleitung gewünscht, die Oberfläche einzelner Länder, wenn
